%! AP Statistics — Unit 1, Lesson 10
\documentclass[10pt]{article}
\usepackage[margin=0.6in]{geometry}
\usepackage{xcolor}
\usepackage[most]{tcolorbox}
\usepackage{enumitem}
\usepackage{multicol}
\usepackage{amsmath,amssymb}
\usepackage{array}
\usepackage{tabularx}
\tcbuselibrary{breakable}

\definecolor{sky}{HTML}{EAF3FB}
\definecolor{navy}{HTML}{1F3A5F}
\definecolor{redbox}{HTML}{FCECEC}
\definecolor{greenbox}{HTML}{EDF8F0}
\definecolor{goldbox}{HTML}{FFF7E6}
\newtcolorbox{skillbox}[2][]{
    colback=#2, colframe=navy, boxrule=0.8pt, arc=2mm,
    left=2mm,right=2mm,top=1.5mm,bottom=1.5mm,
    fonttitle=\bfseries, title={#1}, halign title=center, breakable
}
\setlist[itemize]{leftmargin=1.2em,itemsep=0.35em,topsep=0.35em}
\setlist[enumerate]{leftmargin=1.4em,itemsep=0.35em,topsep=0.35em}
\newcolumntype{C}[1]{>{\centering\arraybackslash}p{#1}}
\newcolumntype{Y}{>{\centering\arraybackslash}X}

\newcommand{\CourseName}{AP Statistics}
\newcommand{\SchoolYear}{2026--2027}
\newcommand{\MeetingLength}{55 minutes}
\newcommand{\UnitNumberName}{Unit 1: Exploring One-Variable Data \quad}
\newcommand{\LessonNumberName}{Lesson 1.10: The Normal Distribution}

\begin{document}
\begin{center}
    {\Large\bfseries \CourseName: \SchoolYear\ (\MeetingLength) \\
    {\small\bfseries \UnitNumberName \LessonNumberName}}
\end{center}

\begin{tcolorbox}[colback=sky,colframe=navy,boxrule=0.9pt,arc=2mm,
    left=2mm,right=2mm,top=1.5mm,bottom=1.5mm]
    \textbf{Primary Objective:} Students will describe the properties of the Normal distribution,
    apply the 68--95--99.7 (Empirical) Rule, calculate and interpret $z$-scores, and use Normal
    distribution tables or technology to find probabilities and percentiles. Students will assess
    whether a dataset is approximately Normal (Big Idea: VAR; AP Skills 2 \& 3).
\end{tcolorbox}

\begin{skillbox}[Priority Ideas \& Skills]{goldbox}
\begin{minipage}[t]{0.48\linewidth}
    \textbf{AP Skills 2 \& 3 --- Data Analysis \& Using Probability and Simulation}
    \begin{itemize}
        \item Apply the 68--95--99.7 Rule to find approximate proportions within 1, 2, or 3 standard deviations of the mean.
        \item Calculate a $z$-score and interpret it in context.
        \item Use Normal distribution tables or technology to find probabilities and percentiles.
        \item Assess Normality using a histogram, dotplot, or Normal probability plot.
    \end{itemize}
\end{minipage}
\hfill
\begin{minipage}[t]{0.48\linewidth}
    \textbf{Key Understandings}
    \begin{itemize}
        \item The Normal distribution is fully described by $\mu$ and $\sigma$; written $N(\mu, \sigma)$.
        \item A $z$-score standardizes a value: how many standard deviations above or below the mean it falls.
        \item Many real distributions are approximately Normal, allowing probability calculations.
        \item The Normal model is an idealization --- real data should be checked before applying it.
    \end{itemize}
\end{minipage}
\end{skillbox}

\begin{skillbox}[{Vocabulary, Concepts, \& Theorems}]{greenbox}
\begin{minipage}[t]{0.48\linewidth}
    \begin{itemize}
        \item \textbf{Normal distribution}: symmetric, bell-shaped; described by mean $\mu$ and standard deviation $\sigma$; notation $N(\mu, \sigma)$
        \item \textbf{68--95--99.7 Rule (Empirical Rule)}: approximately 68\%, 95\%, 99.7\% of data fall within $1\sigma$, $2\sigma$, $3\sigma$ of the mean
        \item \textbf{$z$-score}: $z = \dfrac{x - \mu}{\sigma}$; number of standard deviations a value is from the mean
        \item \textbf{Standard Normal distribution}: $N(0, 1)$; all $z$-scores follow this distribution
    \end{itemize}
\end{minipage}
\hfill
\begin{minipage}[t]{0.48\linewidth}
    \begin{itemize}
        \item \textbf{Normal table (Table A / $z$-table)}: gives the area to the \textit{left} of a given $z$-score
        \item \textbf{Percentile}: the value at or below which a given proportion of the distribution falls
        \item \textbf{Normal probability plot (NPP)}: if points fall close to a diagonal line, the distribution is approximately Normal
        \item \textbf{Assessing Normality}: check histogram (bell shape?), NPP (roughly linear?), and consider context
    \end{itemize}
\end{minipage}
\end{skillbox}

\begin{skillbox}[Activate Prior Knowledge \& Spiral Review]{sky}
\begin{minipage}[t]{0.48\linewidth}
    \textbf{Quick Review (3 min)}\\
    Ask: ``What does a symmetric, bell-shaped histogram look like?'' Students sketch from memory. Then: ``If the mean is 70 and the SD is 10, where does most of the data fall?'' This motivates the Empirical Rule.
\end{minipage}
\hfill
\begin{minipage}[t]{0.48\linewidth}
    \textbf{Spiral Connections}
    \begin{itemize}
        \item Mean and standard deviation (Lesson 1.7) are the two parameters of the Normal model.
        \item Lesson 1.6: symmetric bell shape is one of the named distribution shapes.
        \item Unit 5: the Central Limit Theorem produces Normal sampling distributions --- the Normal model is the backbone of all inference.
    \end{itemize}
\end{minipage}
\end{skillbox}

\begin{skillbox}[Hook \& Lesson]{sky}
\begin{minipage}[t]{0.48\linewidth}
    \textbf{Hook (5 min)}\\
    Display a histogram of heights for 1{,}000 adult men. Ask: ``What percentage are between 67 and 73 inches if the mean is 70 in.\ and SD is 3 in.?'' Students guess. Reveal: the Empirical Rule says $\approx 68\%$. This motivates the power of knowing the distribution is Normal.
\end{minipage}
\hfill
\begin{minipage}[t]{0.48\linewidth}
    \textbf{Lesson Flow (18 min)}
    \begin{enumerate}
        \item State and apply the 68--95--99.7 Rule with the height example.
        \item Introduce $z$-scores; compute and interpret 2--3 examples (positive, negative, $z=0$).
        \item Model using the Normal table: find $P(Z < z)$, $P(Z > z)$, and $P(z_1 < Z < z_2)$.
        \item Find a value given a percentile (inverse Normal).
        \item Show a Normal probability plot and explain how to assess Normality visually.
    \end{enumerate}
\end{minipage}
\end{skillbox}

\begin{skillbox}[Explicit Instruction \& Active Monitoring]{sky}
\begin{minipage}[t]{0.48\linewidth}
    \textbf{Direct Instruction (10 min)}
    \begin{itemize}
        \item Normal table drill: ``The table gives area to the \textit{left}. To find area to the right, subtract from 1. To find area between two $z$-scores, subtract the smaller area from the larger.''
        \item Calculator commands: \texttt{normalcdf(low, high, $\mu$, $\sigma$)} and \texttt{invNorm(area, $\mu$, $\sigma$)}.
        \item Interpretation template: \textit{``A value of $x = 85$ has a $z$-score of $+1.5$, meaning it is 1.5 standard deviations above the mean.''}
    \end{itemize}
\end{minipage}
\hfill
\begin{minipage}[t]{0.48\linewidth}
    \textbf{Active Monitoring}
    \begin{itemize}
        \item Common errors: forgetting to subtract from 1 for ``greater than'' problems; using raw $x$ values in the $z$-table instead of $z$-scores.
        \item Probe: ``Is your answer reasonable? If the $z$-score is negative, should the probability be less than or greater than 0.50?''
        \item Ask students to draw and shade the Normal curve for each problem --- reduces errors.
    \end{itemize}
\end{minipage}
\end{skillbox}

\begin{skillbox}[Group Work \& Differentiation]{redbox}
\begin{minipage}[t]{0.48\linewidth}
    \textbf{Group Activity (8--10 min)}\\
    Context: SAT scores are approximately $N(1060, 195)$. Groups find:
    \begin{enumerate}
        \item $P(\text{score} > 1255)$ using the Empirical Rule.
        \item $P(\text{score} < 900)$ using $z$-scores and the Normal table.
        \item The score at the 90th percentile (inverse Normal).
        \item Whether a score of 1500 is unusually high (justify with $z$-score).
    \end{enumerate}
\end{minipage}
\hfill
\begin{minipage}[t]{0.48\linewidth}
    \textbf{Differentiation}
    \begin{itemize}
        \item \textit{Support}: Provide a Normal curve diagram to shade before computing; Normal table with highlighted rows.
        \item \textit{Extension}: Students investigate what happens when the distribution is \textit{not} Normal (skewed data) --- can we still use the Normal table? Preview of conditions for inference.
        \item \textit{Technology}: Compare table answers with calculator \texttt{normalcdf} outputs.
    \end{itemize}
\end{minipage}
\end{skillbox}

\begin{skillbox}[Individual Work \& Assessment]{redbox}
\begin{minipage}[t]{0.48\linewidth}
    \textbf{Exit Ticket (5 min)}\\
    Assume adult female heights $\sim N(65, 2.5)$ inches. Find: (1) the $z$-score for a woman who is 70 inches tall; (2) the proportion of women taller than 70 inches; (3) the height at the 25th percentile.
\end{minipage}
\hfill
\begin{minipage}[t]{0.48\linewidth}
    \textbf{AP-Style Check}\\
    \textit{``The weights of newborns are approximately Normally distributed with mean 7.5 lb and standard deviation 1.1 lb. What proportion of newborns weigh between 6.4 and 8.6 pounds? Show your work and interpret your answer in context.''}\\[4pt]
    Assess: $z$-scores computed correctly, Normal table or technology used, probability stated and interpreted with context.
\end{minipage}
\end{skillbox}

\begin{skillbox}[Reinforcement \& Extension]{goldbox}
\begin{minipage}[t]{0.48\linewidth}
    \textbf{Homework / Practice}
    \begin{itemize}
        \item Normal distribution problem set: find probabilities, find values given percentiles, interpret $z$-scores in context.
        \item Personal Progress Check 1 (AP Classroom) assigned for Unit 1 review.
    \end{itemize}
\end{minipage}
\hfill
\begin{minipage}[t]{0.48\linewidth}
    \textbf{Extension / Preview}
    \begin{itemize}
        \item \textit{Extension}: Collect class data, compute mean and SD, assess Normality with a histogram. Does the Normal model fit? What would make it more or less appropriate?
        \item \textit{Unit 2 Preview}: Next unit, we explore relationships between \textit{two} variables. How might the Normal model appear when we study two quantitative variables together?
    \end{itemize}
\end{minipage}
\end{skillbox}

\end{document}
